%% Chapter 1: general description of the work.
%% Sections 1.1/1.3 intentionally always show one Spanish and one English
%% blurb (Resumen/Abstract), independent of \thesislanguage, since Spanish
%% universities typically require both regardless of the thesis language.
\chapter{\bilingual{Descripción general del trabajo}{General Description of the Work}}

\section{Resumen}

Resumen extendido del trabajo: contexto, problema, objetivos, solución
propuesta, tecnologías empleadas y principales resultados.
% Demo of the draft-review markers from extracommands.sty -- delete these
% three calls once the real content above is written (or keep sprinkling
% your own \TODO{}/\augusto{}/\tutortwo{}/\student{} while drafting; see
% README.md's "Comentarios de revisión y TODOs").
\TODO{ampliar este resumen con el contenido real del trabajo}%
\augusto{¿puedes precisar más la motivación en el primer párrafo?}%
\student{lo actualizo en la próxima revisión}

\section{Palabras clave}

Palabra clave 1, palabra clave 2, palabra clave 3, palabra clave 4.
\newpage
\section{Abstract}

\begin{otherlanguage}{english}
Extended abstract of the thesis: context, problem, objectives, proposed
solution, technologies used and main results.
\end{otherlanguage}

\section{Keywords}

\begin{otherlanguage}{english}
Keyword 1, keyword 2, keyword 3, keyword 4.
\end{otherlanguage}
